% TOP_PROBLEM: {"id":30007592,"problem_number":"LOCAL-30007592","title":"Risk-premium foundations","category_id":21,"category":{"id":21,"name":"economics","display_name":"Economics","description":"Open economic theory statements, published research agendas, and proposed scoped specifications; question provenance and historical priority are qualified per record.","slug":"economics","order_index":21,"created_at":"2026-10-08T00:00:00Z"},"status":"research_question","external_url":null,"legacy_ids":[],"published":false,"statement":"In the explicitly specified setting: Which empirically plausible combination of preferences, disaster risk, heterogeneity and market frictions explains risk premia across equities and housing? The mathematical statement above fixes the target, assumptions and scope of this catalogue formulation.","economics":{"original_id":81,"field":"Asset pricing and financial markets","kind":"theoretical","formalization_basis":"adapted_research_question","source_status":"explicit_open","priority_status":"earliest_found","priority_note":"The 1985 paper originates the narrower calibrated equity-premium challenge. Housing and broader mechanisms are later additions; original priority for the combined question is not established.","earliest_verified_reference":{"authors":["Rajnish Mehra","Edward C. Prescott"],"title":"The Equity Premium: A Puzzle","year":1985,"url":"https://rajnishmehra.com/wp-content/uploads/2023/09/The-Equity-Premium-A-Puzzle.pdf","doi":"10.1016/0304-3932(85)90061-3","locator":"Section 1, pp.145–146; Section 5, pp.158–159","evidence_summary":"Explicitly asks whether equilibrium models explain the equity premium."},"current_reference":{"authors":["Òscar Jordà","Moritz Schularick","Alan M. Taylor"],"title":"The Total Risk Premium Puzzle","year":2019,"url":"https://www.frbsf.org/wp-content/uploads/wp2019-10.pdf","doi":"10.24148/wp2019-10","locator":"Abstract; Section 7, printed p.40 (PDF p.43)","evidence_summary":"Joint housing-equity premium challenge persists in examined models."},"formalization_note":"Extends a verified equity-premium puzzle to a joint housing-equity identification task. Many model-specific explanations exist; no universal impossibility or current unsolvability theorem is claimed."},"statusQualification":"Published question or research agenda verified at the cited locator. The mathematical specification is adapted for this catalogue and includes additional assumptions; source publication does not establish first-ever priority or certify this exact model remains unsolved. Extends a verified equity-premium puzzle to a joint housing-equity identification task. Many model-specific explanations exist; no universal impossibility or current unsolvability theorem is claimed.","statusReviewedAt":"2026-10-08"}
% FIRST_POSTED: 2026-10-08
% ENTRY_KIND: statement_only
% !TeX program = lualatex
\documentclass[11pt]{article}
\usepackage{fontspec}
\usepackage[a4paper,margin=25mm]{geometry}
\usepackage{amsmath,amssymb,mathtools}
\usepackage{xurl}
\usepackage[hidelinks]{hyperref}
\newcommand{\sref}[2]{\hyperlink{source:#1:#2}{[S#2]}}
\setlength{\emergencystretch}{3em}
\begin{document}
% problemId:30007592
\section*{Risk-premium foundations}\label{30007592}
\subsection{Definitions and mathematical statement}
Fix a discrete-time exchange economy with countries, households and assets. Household \(i\) consumes \(c_{it}>0\), receives labor income, and trades equity, housing claims and a one-period safe bond subject to specified borrowing and trading constraints. Let \(R^e_{t+1}\), \(R^h_{t+1}\), and \(R^f_t\) denote their real gross returns; housing returns include net rents. A model specifies preferences, a probability law for consumption and dividends, beliefs and feasible trading sets. An equilibrium requires household optimality, market clearing and a positive stochastic discount factor \(M_{t+1}\) satisfying the relevant marginal-investor pricing equations. The representative-agent benchmark imposes
\[M_{t+1}=\beta(c_{t+1}/c_t)^{-\gamma},\qquad E_t[M_{t+1}R^j_{t+1}]=1.\]
Here \(E_t\) means expectation conditional on date-\(t\) information, \(\beta\in(0,1)\) is the time-discount factor, \(\gamma>0\) is relative risk aversion, and \(j\) indexes unconstrained traded claims. The research task is to find a quantitatively disciplined extension that jointly reproduces equity and housing excess-return means, their variances, cross-covariance and the safe rate, using preference and friction parameters independently compatible with household evidence. Fix the target population, measurement conventions and acceptable moment tolerances before calibration. Report which restrictions fail in the benchmark, which mechanism relaxes them, and whether distinct mechanisms remain observationally equivalent. Existence of a fitted model is not uniqueness of its explanation.

\par\textbf{Formulation and scope.} Extends a verified equity-premium puzzle to a joint housing-equity identification task. Many model-specific explanations exist; no universal impossibility or current unsolvability theorem is claimed.

\par\textbf{Open-question provenance.} An explicit question or research agenda was verified at the source locator. This does not by itself certify that every later version remains unresolved \sref{30007592}{1}.

\par\textbf{Historical priority.} The 1985 paper originates the narrower calibrated equity-premium challenge. Housing and broader mechanisms are later additions; original priority for the combined question is not established.

\subsection{Short English statement}
In the explicitly specified setting: Which empirically plausible combination of preferences, disaster risk, heterogeneity and market frictions explains risk premia across equities and housing? The mathematical statement above fixes the target, assumptions and scope of this catalogue formulation.

\subsection{Sources}
\begin{itemize}\raggedright
\item[\textnormal{[S1]}]\hypertarget{source:30007592:1}{} Rajnish Mehra, Edward C. Prescott. 1985. The Equity Premium: A Puzzle. Section 1, pp.145–146; Section 5, pp.158–159.
\url{https://rajnishmehra.com/wp-content/uploads/2023/09/The-Equity-Premium-A-Puzzle.pdf}.
DOI: 10.1016/0304-3932(85)90061-3.
{\small\textit{Earliest verified question/agenda reference; historical priority qualified below. Explicitly asks whether equilibrium models explain the equity premium.}}

\item[\textnormal{[S2]}]\hypertarget{source:30007592:2}{} Òscar Jordà, Moritz Schularick, Alan M. Taylor. 2019. The Total Risk Premium Puzzle. Abstract; Section 7, printed p.40 (PDF p.43).
\url{https://www.frbsf.org/wp-content/uploads/wp2019-10.pdf}.
DOI: 10.24148/wp2019-10.
{\small\textit{Supporting or current literature; see evidence qualification. Joint housing-equity premium challenge persists in examined models.}}
\end{itemize}
\end{document}
